\documentclass[11pt]{article}

% Change "review" to "final" to generate the final (sometimes called camera-ready) version.
% Change to "preprint" to generate a non-anonymous version with page numbers.
\usepackage[preprint]{acl}

\usepackage{graphicx}
\usepackage{booktabs}
% Standard package includes
\usepackage{times}
\usepackage{latexsym}

% For proper rendering and hyphenation of words containing Latin characters (including in bib files)
\usepackage[T1]{fontenc}
% For Vietnamese characters
% \usepackage[T5]{fontenc}
% See https://www.latex-project.org/help/documentation/encguide.pdf for other character sets

% This assumes your files are encoded as UTF8
\usepackage[utf8]{inputenc}

% This is not strictly necessary, and may be commented out,
% but it will improve the layout of the manuscript,
% and will typically save some space.
\usepackage{microtype}

% This is also not strictly necessary, and may be commented out.
% However, it will improve the aesthetics of text in
% the typewriter font.
\usepackage{inconsolata}

%Including images in your LaTeX document requires adding
%additional package(s)
\usepackage{graphicx}
\usepackage{amsmath}

% If the title and author information does not fit in the area allocated, uncomment the following
%
%\setlength\titlebox{<dim>}
%
% and set <dim> to something 5cm or larger.

\title{The Robots Are Coming (Again): Job-Loss Fear in Online Discussions before and after ChatGPT}

% Author information can be set in various styles:
% For several authors from the same institution:
% \author{Author 1 \and ... \and Author n \\
%         Address line \\ ... \\ Address line}
% if the names do not fit well on one line use
%         Author 1 \\ {\bf Author 2} \\ ... \\ {\bf Author n} \\
% For authors from different institutions:
% \author{Author 1 \\ Address line \\  ... \\ Address line
%         \And  ... \And
%         Author n \\ Address line \\ ... \\ Address line}
% To start a separate ``row'' of authors use \AND, as in
% \author{Author 1 \\ Address line \\  ... \\ Address line
%         \AND
%         Author 2 \\ Address line \\ ... \\ Address line \And
%         Author 3 \\ Address line \\ ... \\ Address line}

\author{
  \textbf{Michał Bortkiewicz},
  \textbf{Aleksandra Czarniecka},
  \textbf{Mikołaj Drojma},
  \textbf{Zuzanna Nasiłowska},
  \textbf{Patrycja Smits}  
  \\
  Faculty of Information and Communication Technology, \\
  Wrocław University of Science and Technology
}
%\author{
%  \textbf{First Author\textsuperscript{1}},
%  \textbf{Second Author\textsuperscript{1,2}},
%  \textbf{Third T. Author\textsuperscript{1}},
%  \textbf{Fourth Author\textsuperscript{1}},
%\\
%  \textbf{Fifth Author\textsuperscript{1,2}},
%  \textbf{Sixth Author\textsuperscript{1}},
%  \textbf{Seventh Author\textsuperscript{1}},
%  \textbf{Eighth Author \textsuperscript{1,2,3,4}},
%\\
%  \textbf{Ninth Author\textsuperscript{1}},
%  \textbf{Tenth Author\textsuperscript{1}},
%  \textbf{Eleventh E. Author\textsuperscript{1,2,3,4,5}},
%  \textbf{Twelfth Author\textsuperscript{1}},
%\\
%  \textbf{Thirteenth Author\textsuperscript{3}},
%  \textbf{Fourteenth F. Author\textsuperscript{2,4}},
%  \textbf{Fifteenth Author\textsuperscript{1}},
%  \textbf{Sixteenth Author\textsuperscript{1}},
%\\
%  \textbf{Seventeenth S. Author\textsuperscript{4,5}},
%  \textbf{Eighteenth Author\textsuperscript{3,4}},
%  \textbf{Nineteenth N. Author\textsuperscript{2,5}},
%  \textbf{Twentieth Author\textsuperscript{1}}
%\\
%\\
%  \textsuperscript{1}Affiliation 1,
%  \textsuperscript{2}Affiliation 2,
%  \textsuperscript{3}Affiliation 3,
%  \textsuperscript{4}Affiliation 4,
%  \textsuperscript{5}Affiliation 5
%\\
%  \small{
%    \textbf{Correspondence:} \href{mailto:email@domain}{email@domain}
%  }
%}

\begin{document}
\maketitle
\begin{abstract}
TBA
\end{abstract}

\section{Introduction}

Fears that machines will take people's jobs are not new. Nearly a century ago, \citet{keynes1931} warned of \textit{technological unemployment}, and every other wave of automation has revived similar concerns \citep{mokyr2015}.
About a decade ago, \citet{frey2017} estimated that 47\% of U.S. jobs were at high risk of computerisation, and in 2017, five years before ChatGPT, 72\% of EU citizens already agreed that robots \textit{steal people's jobs} \citep{eurobarometer2017}.

The launch of ChatGPT on 30 November 2022 brought the robots back into the spotlight, this time as chatbots capable of generating text, writing code, automating repetitive tasks and more.

Since then, AI-powered tools have spread quickly into everyone's everyday work and study. Their impact on employment and the labor market has attracted increasing attention. Estimations suggest that many jobs may be affected by this transformation. In Poland, for instance, between 22\% and 30\% of workers are employed in occupations that are exposed to generative AI \citep{pie2024,ilonask2025}. Public concerns are widespread: 52\% of U.S. workers are worried about the future use of AI at work, and 32\% think it will lead to fewer job opportunities for them \citep{pew2025}; 66\% of EU citizens agree that AI and robots \textit{steal people's jobs} \citep{eurobarometer2025}. At the same time, AI is often presented as a technology that can increase productivity and support workers, especially less experienced ones \citep{brynjolfsson2025qje}. Many people hold both views at once: 66\% of EU workers see the impact of recent digital technologies, including AI, on their current job as positive, and 63\% agree that robots and AI are good for society because they help people do their jobs \citep{eurobarometer2025}. This ambivalence raises an important question.

After the launch of ChatGPT, did people become more afraid for their jobs and their future, or did they start to see AI as an opportunity that makes their lifes and work easier?

Existing surveys cannot answer this well. They aren't conducted regularly and often depend on whether respondents are asked about the labor market in general or their own jobs. Online discussions, on the other hand, provide a continuous record of public opinion expressed in people's own words.

We study this question on two platforms in two languages: Wykop.pl, a Polish microblogging site (about 125,000 posts), and YouTube comments under videos about work and technology, mostly in English (about 50,000 comments). We hypothesize that the number of posts expressing concerns about job loss in connection with AI increases after November 2022 (H1). Furthermore, we hypothesize that, after ChatGPT launch, job advertisements contain a larger number of skill requirements (H2). To test these hypotheses, we developed a three-stage pipeline (\S\ref{sec:annotation}) that enabled us to identify relevant discussions and analyze the emotional tone and attitudes expressed toward AI and its impact on work.

\section{Related Work}
Previous research has examined how hopes and concerns about AI are expressed in public discourse. Analysing three decades of New York Times coverage, \citet{fast2017} found that generally optimistic reporting coexisted with growing concerns about AI’s effects on work. More directly related to our study, \citet{fodor2026} analysed over ten years of Reddit discussions about AI and the labour market. They reported mixed attitudes and a more positive outlook following the release of ChatGPT, including a shift towards discussing AI as a tool for augmenting human work. These findings suggest that increased attention to AI does not necessarily imply increased fear of job loss. Meanwhile, \citet{xu2024} used topic modelling and sentiment analysis to study Reddit discussions of ChatGPT, identifying broadly favourable attitudes alongside discussion of its societal implications. However, general sentiment towards a technology does not directly measure employment-related fear. Our study therefore focuses specifically on expressions of job-loss concern, comparing discussions before and after ChatGPT’s release on Wykop and YouTube. This extends the platform and language coverage of the studies discussed above, while treating the observed texts as evidence about discourse within the sampled communities. Methodologically, our approach follows \citet{grimmer2013}, who emphasise that automated text analysis requires validation tailored to the substantive concept being measured.




\section{Data}

We collected data from two platforms that differ in language and audience (Table~\ref{tab:corpus}).

Wykop.pl is one of the largest Polish microblogging sites. Its users discuss a broad range of topics, including technology, work, politics and current events. We collected about 125,000 posts and comments published between January 2019 and October 2026. We retrieved all posts containing at least one of eight tags related to work, IT and AI: \#praca ("work"), \#programista15k ("15k programmer", a popular tag referring to software developers and IT careers), \#pracait ("IT jobs"), \#it ("information technology"), \#rekrutacja ("recruitment"), \#ai ("artificial intelligence"), \#chatgpt and \#sztucznainteligencja ("artificial intelligence").

To cover English-language discussion, we collected about 50,000 comments under YouTube videos about work and technology, retrieved with nine search phrases (Appendix~\ref{app:queries}). We date each text by the time the comment was posted, and the corpus spans May 2019 to October 2026. The corpus is restricted to English. (details in Appendix~\ref{app:queries}).

First, we preprocessed the data by removing duplicates, posts with only links, images or hashtags, and automatically generated content, such as the headers and footers. Usernames were replaced with salted SHA-256 hashes, and e-mail addresses, phone numbers and user mentions were masked in the text. Next, we divided the dataset into two periods based on the release of ChatGPT on 30 November 2022. The \textit{before} period covers all content published up to 29 November 2022 and the \textit{after} period begins on 30 November 2022. For the time-series analysis, texts were aggregated by calendar month.

\begin{table}[t]
\centering
\small
\begin{tabular}{lrrrr}
\hline
 & \multicolumn{2}{c}{\textbf{Wykop}} & \multicolumn{2}{c}{\textbf{YouTube}} \\
 & before & after & before & after \\
\hline
Texts            & 59,419 & 60,633 & [--] & [--] \\
\% work          & 74.4   & 63.2   & [--] & [--] \\
\% AI            & 2.9    & 26.8   & [--] & [--] \\
\% AI $\wedge$ work & 0.9 & 7.6    & [--] & [--] \\
\hline
\end{tabular}
\caption{Corpus statistics before (until 30 Nov 2022) and after (from 30 Nov 2022) the release of ChatGPT. The work and AI tags come from the stage-1 filter (\S\ref{sec:annotation}).}
\label{tab:corpus}
\end{table}


\section{Annotation \& Models}
\label{sec:annotation}

\subsection{Annotation scheme}

The unit of annotation is an individual post or comment.
We classify only the author's text, without providing the model
with publication dates, video titles, linked content, or parent
posts. The annotation guidelines distinguish explicit statements
from opinions inferred from context. Hashtags alone are not
sufficient evidence of relevance, and quoted concerns are not
treated as the author's position unless endorsed. Irony is
interpreted according to its intended meaning where this can
be established from the text.

We first assign two independent binary labels:
\texttt{ai}, indicating a reference to artificial intelligence,
and \texttt{work}, indicating discussion of employment, careers,
recruitment, wages, or professional tasks and skills.
Generic references to software, algorithms, or automation
without a clear AI reference receive \texttt{ai=no}.

All texts with \texttt{work=yes}, whether or not they mention AI,
receive three further labels, which allow the pre- and
post-ChatGPT periods to be compared.
\texttt{job\_insecurity} records whether the author expresses
worry about their own or others' job security, income, or
ability to find work, regardless of its source.
\texttt{insecurity\_cause}, assigned only when insecurity is
expressed, identifies the cited source: AI or automation,
economic conditions and layoffs, offshoring or immigration,
employer practices, individual factors such as age or skills,
other, or no cause stated.
\texttt{sector} identifies the principal occupational domain.

For texts discussing both AI and work, \texttt{stance} captures
the author's assessment of AI's employment implications:
\texttt{threat}, \texttt{reassurance}, \texttt{opportunity},
\texttt{mixed}, or \texttt{none}. Threat includes anticipated
job loss, reduced hiring, lower wages, and occupational or
skill obsolescence. Mixed statements combine threats with
opportunities or reassurance. Merely mentioning AI use at work
does not establish a stance.

Detailed annotation of these texts adds three dimensions.
\texttt{attitude} records a positive, negative, mixed, or neutral
evaluation of AI in the labour market.
\texttt{threat\_type}, assigned only to threat or mixed
statements, distinguishes job loss, hiring difficulties,
lower pay, skill obsolescence, and other concerns.
\texttt{competence\_claim} records whether the author claims
that AI increases, reduces, or leaves skill requirements
unchanged, with a separate category for no such claim.
These labels describe expressed views rather than
established effects on employment.

\subsection{Model pipeline}

Both platforms follow the same three steps with the same label scheme and
instructions: a screening pass over the whole corpus, a double annotation of
a stratified sample with adjudication of disagreements, and labelling of the
whole corpus with a model validated on the adjudicated sample.

First, a screening pass with
GPT-6 Luna assigns \texttt{ai} and \texttt{work} to every text in the
corpus, with responses constrained to a JSON schema.

Second, for each platform we selected 8,000 texts for a model annotation
pass. The sample is stratified by the screening labels and by period, and
sampling probabilities are retained by stratum and period.

Two coding agents, Claude Opus 5.5 (claude-opus-5-5) and Codex (Sol 6.1, medium reasoning effort), each independently annotated the entire sample
in 200 batches of 40 texts, using the same instructions. Both agents
received the text and identifiers without the screening predictions or
period labels, and reassigned the screening labels before adding the
detailed labels. We then identified disagreements between the two sets
of annotations and, acting as human adjudicators, reviewed these cases
and corrected the annotations to produce a final adjudicated dataset for
each platform. These datasets were subsequently used to train a BERT
classifier.

The two annotators disagreed on 1,424 labels in 1,041 texts. For each disagreement, the authors decided for final labels. The result is a gold set of 8000 texts.


\begin{table}[t]
\centering
\small
\begin{tabular}{lrr}
\hline
Label & Agreement & $\kappa$ \\
\hline
\texttt{ai}               & 0.93 & 0.83 \\
\texttt{work}             & 0.84 & 0.66 \\
\texttt{job\_insecurity}   & 0.90 & 0.70 \\
\texttt{insecurity\_cause} & 0.87 & 0.65 \\
\texttt{sector}           & 0.84 & 0.75 \\
\texttt{stance}           & 0.94 & 0.78 \\
\texttt{attitude}         & 0.89 & 0.76 \\
\texttt{threat\_type}      & 0.96 & 0.78 \\
\texttt{competence\_claim} & 0.98 & 0.64 \\
\hline
\end{tabular}
\caption{Agreement between Claude Opus 5.5 and GPT-6.1 Sol on the Wykop
sample before adjudication. Detailed labels are compared on the texts both
annotators labelled \texttt{work=yes}. The low $\kappa$ of
\texttt{competence\_claim} reflects its rarity (98\% \textit{none}).}
\label{tab:agents}
\end{table}

Since stage 1 already covers the
whole corpus, only the seven stage 2 labels remain to be assigned, and only to
texts with \texttt{work=yes}. We first planned to fine-tune HerBERT
(\texttt{allegro/herbert-base-cased}) on the adjudicated sample, with one
shared encoder and a classification head per label, each head trained only on
the texts where its label applies. The work-related gold texts were split
into training, validation and test sets (70/15/15\%, stratified by period and
stratum. The labels, however, are strongly
imbalanced: \texttt{insecurity\_cause}, \texttt{threat\_type} and
\texttt{competence\_claim} apply to only 207 to 622 training texts, spread over
up to seven classes, several of them with only a handful of examples. Even
with class-weighted losses, the model overfitted the training set and did not
generalise to the validation and test sets, so we did not use it.  Instead,
the corpus was labelled by GPT-6 Luna (low reasoning effort),
with the stage 2 instruction and responses constrained to a JSON schema. On
the 473 adjudicated test texts it reached a mean macro-F1 of 0.67, computed
for each label on the texts where the gold label applies. 

\subsection{Quality control and validation}

Model outputs are checked for allowed label values and
conditional consistency. The review importer also
checks text identifiers, ordering, and completeness of each batch.
These checks establish output integrity, not semantic accuracy.

For Wykop, the screening pass is validated on a reference sample of 150
texts, drawn in equal parts from texts an earlier classifier had labelled
as both AI- and work-related and from all other texts. Two annotators
labeled the sample independently, following the guidelines. Their
agreement was high for \texttt{ai} (Cohen's $\kappa = 0.86$) and moderate
for \texttt{work} ($\kappa = 0.55$), the label that requires the most
judgement. We treat as gold the labels on which both annotators agree.
Texts consisting only of hashtags were removed in preprocessing and are
not evaluated.

Table~\ref{tab:validation} compares the screening labels with the gold
standard. The classifier agreed with the gold labels more closely than the
two annotators agreed with each other: $\kappa = 0.98$ for \texttt{ai},
$0.80$ for \texttt{work} and $0.92$ for both labels together. It also agreed
with each annotator separately at least as well as the annotators agreed
with each other (\texttt{ai}: $\kappa = 0.94$ and $0.91$; \texttt{work}:
$0.70$ and $0.58$). Most errors concern \texttt{work}, the label on which
the annotators also disagreed most. 

\begin{table}[t]
\centering
\small
\begin{tabular}{lrrrrr}
\hline
Label & $\kappa_{\text{human}}$ & $n$ & $\kappa_{\text{model}}$ & Acc. & P / R \\
\hline
\texttt{ai}              & 0.86 & 134 & 0.98 & 99.3 & 1.00 / 0.98 \\
\texttt{work}            & 0.55 & 119 & 0.80 & 93.3 & 0.96 / 0.96 \\
\texttt{ai} $\wedge$ \texttt{work} & 0.72 & 125 & 0.92 & 96.8 & 1.00 / 0.88 \\
\hline
\end{tabular}
\caption{Screening-pass validation on Wykop. $\kappa_{\text{human}}$:
agreement between the two annotators; $n$: texts with an agreed (gold)
label; $\kappa_{\text{model}}$, accuracy (\%), precision and recall (for
\textit{yes}): the screening classifier against the gold label. 95\%
bootstrap intervals for $\kappa_{\text{model}}$: 0.95--1.00, 0.64--0.91,
0.83--0.98.}
\label{tab:validation}
\end{table}

\section{Results}
\label{sec:results}

We test H1 with three indicators: the share of work-related texts expressing
job insecurity, the share of insecure texts that name AI as the cause (Wykop) or
also discuss AI (YouTube), and the stance towards AI's effects on employment.
All comparisons are between the period before (until 29~November~2022) and
after (from 30~November~2022) the release of ChatGPT. Because Wykop is labelled
in full, its tests use the whole population of texts; percentages for YouTube
are weighted estimates from the adjudicated sample.

\subsection{Job insecurity on Wykop}

Table~\ref{tab:wykop-insecurity} summarises job insecurity among the 82,538
work-related Wykop texts. The share of texts expressing job insecurity rose
from 12.8\% before to 20.4\% after the release of ChatGPT (risk ratio
$RR=1.59$, 95\% CI 1.54--1.64; $\chi^2(1)=854$, $p<0.001$). The composition of
causes changed as well ($\chi^2(6)=1251$, $p<0.001$, Cramér's $V=0.30$). AI
grew from 0.7\% to 14.7\% of all stated causes, and economic conditions from
16.1\% to 21.3\%, while individual circumstances fell from 36.2\% to 19.2\%.
Measured against all work-related texts, insecurity attributed to AI increased
from 0.09\% to 2.99\% ($RR=32$, 95\% CI 24--44).

\begin{table}[t]
\centering\small
\begin{tabular}{lrr}
\toprule
 & before & after \\
\midrule
Work-related texts & 44,203 & 38,335 \\
\quad with job insecurity & 12.8\% & \textbf{20.4\%} \\
\midrule
\multicolumn{3}{l}{\emph{Cause (share of insecure texts)}} \\
AI & 0.7\% & \textbf{14.7\%} \\
Economy / crisis & 16.1\% & 21.3\% \\
Employer & 21.9\% & 16.4\% \\
Individual situation & 36.2\% & 19.2\% \\
Offshoring / immigration & 1.0\% & 2.6\% \\
Other & 8.4\% & 7.4\% \\
No cause stated & 15.7\% & 18.4\% \\
\bottomrule
\end{tabular}
\caption{Job insecurity in work-related Wykop texts before and after the
release of ChatGPT. $\chi^2$ tests: insecurity share $p<0.001$; cause
distribution $p<0.001$, $V=0.30$.}
\label{tab:wykop-insecurity}
\end{table}

\paragraph{Timing.} A simple before--after comparison cannot tell whether
insecurity was already rising before November 2022. We therefore fitted an
interrupted time-series model to monthly shares (January 2019 -- September
2026, 93 months), with a linear trend, a level change and a trend change at
the release date. Months were weighted by the number of texts, and standard
errors are Newey--West (HAC, 6 lags). Before the release, neither overall nor
AI-related insecurity had a trend ($p=0.94$ and $p=0.16$). AI-attributed
insecurity did not jump immediately (level change $+0.5$~pp, $p=0.37$), but its
trend increased by 0.10~pp per month afterwards ($p<0.001$). Overall
insecurity also jumped ($+3.6$~pp, $p<0.001$) and increased faster
($+0.18$~pp per month, $p=0.004$). Insecurity unrelated to AI showed only a
level shift ($+3.0$~pp, $p=0.018$) without a trend change ($p=0.29$).

\begin{figure}[t]
\centering
\includegraphics[width=\columnwidth]{wykop_insecurity_yearly.png}
\caption{Wykop: job insecurity by year, overall and split by whether the text
also discusses AI. Share of work-related texts in each year. *2026: data until
14 October 2026.}
\label{fig:wykop-insecurity}
\end{figure}


\paragraph{Robustness to tag selection.} Wykop texts were collected via tags,
including \#ai, \#chatgpt and \#sztucznainteligencja, which are much more
frequent after the release (0.3\% vs.\ 6.9\% of work-related texts). After
excluding all texts carrying an AI tag, AI-attributed insecurity still rises
from 0.05\% to 1.80\% of work-related texts and overall insecurity from 12.8\%
to 20.1\%. The effect is therefore not an artefact of tag-based collection,
although its size is smaller.

\subsection{Stance towards AI on Wykop}

Among Wykop texts discussing both AI and work (457 before, 4,240 after), the
distribution of stance changed significantly ($\chi^2(4)=195$, $p<0.001$,
$V=0.20$). The share of texts seeing AI as a \emph{threat} to jobs rose from
7.0\% to 28.5\%, but the shares of \emph{opportunity} (1.5\% $\to$ 8.8\%) and
\emph{reassurance} (1.5\% $\to$ 5.1\%) also grew, while texts without an
assessment fell from 89.9\% to 56.2\%. The attitude towards AI in the labour
market became more polarised: negative 8.5\% $\to$ 39.8\%, positive 3.9\%
$\to$ 13.0\%, neutral 87.3\% $\to$ 44.1\% ($\chi^2(3)=309$, $p<0.001$,
$V=0.26$). Among AI-related work texts, job insecurity was expressed in 34.3\%
after the release, compared with 18.6\% of other work-related texts.



\subsection{YouTube}

On YouTube, 36.0\% of the screened comments concern work, 7.3\% concern AI and
3.3\% concern both (weighted estimates). Job insecurity appears in 5.0\% of
work-related comments, considerably less than on Wykop. Among comments
discussing both AI and work, 40.3\% see AI as a threat to jobs, 19.6\% offer
reassurance, 11.1\% see an opportunity and 8.3\% hold a mixed view.

Figures~\ref{fig:yt-stance} and~\ref{fig:yt-insecurity} show the yearly
development. From 2020 to 2022 all stance categories stay below 0.2\% of
comments. From 2023 onward the \emph{threat} share grows every year
(0.6\%, 0.9\%, 1.5\% and 6.6\% in 2026), while reassurance and opportunity grow
more slowly. Job insecurity in AI-related comments is practically absent before
2023 ($\leq 0.04$\%) and rises steadily to 0.95\% in 2026. Overall job insecurity
on YouTube does not follow the same pattern: it fluctuates between 0.7\% and
2.8\% and peaks in 2024, and in 2026 AI-related insecurity for the first time
exceeds insecurity unrelated to AI (0.95\% vs.\ 0.60\%).


\begin{figure}[t]
\centering
\includegraphics[width=\columnwidth]{yt_stance_yearly.png}
\caption{YouTube: stance towards AI's effects on work, by year. Share of all
comments in each year (weighted). *2026: data until 26 October 2026.}
\label{fig:yt-stance}
\end{figure}

\begin{figure}[t]
\centering
\includegraphics[width=\columnwidth]{yt_insecurity_yearly.png}
\caption{YouTube: job insecurity by year, overall and split by whether the
comment also discusses AI. Share of all comments in each year (weighted).}
\label{fig:yt-insecurity}
\end{figure}

\subsection{Content of AI-related concerns on both platforms}

The detailed labels of AI-related work texts are strikingly similar on the two
platforms (Table~\ref{tab:stage2-compare}). When AI is seen as a threat, the
concern is overwhelmingly about job loss (79\% on Wykop, 77\% on YouTube),
far ahead of hiring difficulties, skill obsolescence and lower pay. Claims
about changing skill requirements are rare on both platforms. IT is the most
frequently named sector on both, more so on Wykop, partly because of the
IT-related collection tags.

\begin{table}[t]
\centering\small
\begin{tabular}{lrr}
\toprule
 & Wykop & YouTube \\
\midrule
\multicolumn{3}{l}{\emph{Attitude (AI $\wedge$ work texts)}} \\
Negative & 36.8 & 45.4 \\
Neutral & 48.3 & 35.6 \\
Positive & 12.2 & 15.9 \\
Mixed & 2.7 & 3.1 \\
\midrule
\multicolumn{3}{l}{\emph{Threat type (threat/mixed texts)}} \\
Job loss & 79.2 & 77.0 \\
Hiring & 8.8 & 9.7 \\
Skills obsolescence & 7.3 & 6.0 \\
Pay & 2.5 & 2.9 \\
Other & 2.2 & 4.6 \\
\midrule
\multicolumn{3}{l}{\emph{Competence claim}} \\
None & 94.4 & 87.8 \\
More skills needed & 2.6 & 9.4 \\
Fewer skills needed & 3.0 & 2.2 \\
\midrule
\multicolumn{3}{l}{\emph{Sector}} \\
IT & 55.4 & 38.6 \\
Unspecified & 28.8 & 41.1 \\
\bottomrule
\end{tabular}
\caption{Stage-2 labels of texts discussing AI and work, whole period (\%).
YouTube values are weighted.}
\label{tab:stage2-compare}
\end{table}







\section{Discussion}
\label{sec:discussion}

\paragraph{H1 is supported on both platforms.} In two different communities,
in two languages, job insecurity connected with AI was virtually absent before
November 2022 and grew afterwards. The interrupted time-series analysis of Wykop
adds an important detail: the release of ChatGPT did not produce an immediate
jump, but marked the start of a steady increase that continued through 2026.
The YouTube yearly series shows the same shape. This is consistent with a
gradual diffusion of generative AI into everyday work, as opposed to a single
moment of shock.

\paragraph{AI is not the only source of concern.} On both platforms, job
insecurity unrelated to AI also changed after 2022. On Wykop it rose by about
4~pp, and the largest single increase among causes was economic conditions; on
YouTube it peaked in 2024. The period after ChatGPT coincided with layoffs and
slower hiring, especially in IT, which was the dominant sector in our data. A
before--after design cannot separate these influences, so the growth of overall
insecurity should not be attributed to AI alone. What our data does show is that
AI increasingly became the \emph{reason} people gave for their insecurity: by
2026, AI-related insecurity on YouTube exceeded insecurity unrelated to AI.

\paragraph{More fear, but also more discussion.} The rise of threat statements
was accompanied by a rise in opportunity and reassurance statements. On Wykop,
neutral attitudes gave way to both negative and positive ones. This partly
agrees with \citet{fodor2026gpt}, who found mixed attitudes on Reddit and a
more positive outlook after ChatGPT: on our platforms, negative views grew
faster, but positive views grew too. Increased attention to AI therefore does
not translate one-to-one into fear.

\paragraph{Similar concerns, different levels.} The levels differ
substantially between platforms: job insecurity occurs in about 20\% of
work-related Wykop texts after the release, but in about 5\% of work-related
YouTube comments. This reflects the platforms and the collection methods rather
than countries: Wykop texts come from career and IT tags where users describe
their own situation, while YouTube comments react to a video. The
\emph{content} of AI-related concern, however, is nearly identical: in both
languages, about four in five threat statements concern losing jobs, and claims
about changing skill requirements are rare.


\section{Conclusion}
\label{sec:conclusion}

We compared online discussions of work before and after the release of ChatGPT
on a Polish and an English-language platform. On both, job insecurity related
to AI was nearly absent before November 2022 and grew steadily afterwards,
and threat statements focused on job loss. At the same time, insecurity with
other causes also increased, and opportunity and reassurance statements became
more frequent, so the discourse became more polarised rather than uniformly
fearful. Future work should model the influence of labour market conditions
directly, for example with unemployment or job-vacancy data, and examine
whether the changing discourse is reflected in the skill requirements of job
advertisements.


\section*{Limitations}

This document does not cover the content requirements for ACL or any
other specific venue.  Check the author instructions for
information on
maximum page lengths, the required ``Limitations'' section,
and so on.

\section*{Acknowledgments}


% Bibliography entries for the entire Anthology, followed by custom entries
%\bibliography{custom,anthology-overleaf-1,anthology-overleaf-2}

% Custom bibliography entries only
\bibliography{custom}

\appendix

\section{Search phrases and collection protocol}
\label{app:queries}

Videos were retrieved from the YouTube Data API with the nine search phrases
listed in Table~\ref{tab:queries}. The phrases contain no AI-related terms, so
that the retrieval itself does not favour discussion of AI. They were chosen by
probing YouTube search. We kept phrases that returned on-topic
videos with at least 50 comments and a measurable share of comments about AI
after the release of ChatGPT. Phrases about general workplace themes (``workplace
automation'', ``future of work'', ``job security'') returned mostly webinars and
unrelated videos with few comments, and were dropped.

\begin{table}[h]
\centering
\small
\begin{tabular}{ll}
\hline
Group & Search phrase \\
\hline
Future of jobs and automation & \textit{will robots take our jobs} \\
                              & \textit{jobs that will disappear} \\
                              & \textit{automation unemployment} \\
Jobs exposed to AI            & \textit{is computer science still worth it} \\
                              & \textit{junior developer jobs} \\
                              & \textit{software engineer jobs} \\
                              & \textit{data analyst jobs} \\
General job search (background) & \textit{entry level jobs} \\
                                & \textit{how to get a job} \\
\hline
\end{tabular}
\caption{Search phrases used to retrieve videos.}
\label{tab:queries}
\end{table}

\paragraph{Sampling protocol.}
Searches were run separately for two periods, before ChatGPT
(2019-01-01 to 2022-11-29) and after it (2022-11-30 onward), with results
ordered by relevance, so that both periods are represented equally (target
25{,}000 comments each). A video was kept if its audio or default language was
English, it had at least 10{,}000 views and at least 50 comments. From each
video we collected top-level comments in chronological order. A comment was
kept if it had at least 20 characters and a language detector
(\texttt{langdetect}) rated it English with probability of at least 0.90.
Contact details were masked. The comment's period is determined by its own
posting time, not by the video's upload date.

\end{document}
